\subsection{Binäre Operatoren als Funktionen}

Ein binärer Operator auf einer Menge bestimmt eine Funktion auf dem
kartesischen Quadrat dieser Menge. Die Projektionen geben die beiden
Argumente des Operators an.

\FormulaDefDeltaK[Funktionsgraph eines binären Operators]{%
  \mathsf{BinOp}(A,\star)\vdash
  \begin{aligned}[t]
    &\text{(i)}\;\BinaryOpFunction{A}{\star}\colon A\times A\to A\\[-2pt]
    &\text{(ii)}\;\forall x,y\in A\,
      \bigl(\BinaryOpFunction{A}{\star}(x,y)\coloneqq x\star y\bigr)
  \end{aligned}%
}{BinaryOperationFunctionDef}{%
  \DeltaRow{Mengen}{A\dsep p\dsep x\dsep y}
  \DeltaRow{Funktionen}{\star\dsep G\dsep\BinaryOpFunction{A}{\star}}
  \DeltaRow{\textbf{Neues Symbol}}{}
  \DeltaPrem{Operatorfunktion}{\BinaryOpFunction{A}{\star}}
}
\begin{tabproofsplitwide}
  \proofpartwide{Wohldefiniertheit und Funktionstypisierung}
  \proofstepwidestar[1]{\mathsf{BinOp}(A,\star)}{\rA}
  \proofstepwidestar[2]{p\in A\times A}{\rA}
  \proofstepwidestar[2]{\pi_1(p)\in A}{%
    \FormulaRefAuto{z\in A\times B\vdash\pi_1(z)\in A}[thm]{2}}
  \proofstepwidestar[2]{\pi_2(p)\in A}{%
    \FormulaRefAuto{z\in A\times B\vdash\pi_2(z)\in B}[thm]{2}}
  \proofstepwidestar[1,2]{\pi_1(p)\star\pi_2(p)\in A}{%
    \FormulaRefAuto{BinaryOperationClosureAxiom}[ax]{1,3,4}}
  \proofstepwidestar[1]{%
    \forall p\in A\times A\,
    \pi_1(p)\star\pi_2(p)\in A}{\rUI{\rRI{2,5}}}
  \proofstepwidestar[1]{%
    \begin{aligned}[t]
      &\exists!G\colon A\times A\to A\,\\[-2pt]
      &\quad\forall p\in A\times A\,
        G(p)=\pi_1(p)\star\pi_2(p)
    \end{aligned}}{%
    \FormulaRefAuto{TypedFunctionDefinitionPrinciple}[thm]{6}}
  \proofstepwidestar{%
    \begin{aligned}[t]
      &\BinaryOpFunction{A}{\star}\coloneqq
        \iota G\colon A\times A\to A\,\\[-2pt]
      &\quad\forall p\in A\times A\,
        G(p)=\pi_1(p)\star\pi_2(p)
    \end{aligned}}{\FormulaRefAuto{IotaDefinitionDef}[def]{7}}
  \proofstepwidestar[1]{%
    \begin{aligned}[t]
      &\BinaryOpFunction{A}{\star}\colon A\times A\to A\land{}\\[-2pt]
      &\forall p\in A\times A\,
        \BinaryOpFunction{A}{\star}(p)=\pi_1(p)\star\pi_2(p)
    \end{aligned}}{%
    \FormulaRefAuto{IotaDefinitionDef}[def]{7,8}}
  \proofstepwidestar[1]{%
    \BinaryOpFunction{A}{\star}\colon A\times A\to A}{\rAEa{9}}
  \proofstepwidestar[1]{%
    \forall p\in A\times A\,
    \BinaryOpFunction{A}{\star}(p)
      =\pi_1(p)\star\pi_2(p)}{\rAEb{9}}
  \closeproofpartwide

  \proofpartwide{Auswertung an geordneten Paaren}
  \setcounter{proofstepnr}{11}
  \proofstepwidestar[12]{x\in A}{\rA}
  \proofstepwidestar[13]{y\in A}{\rA}
  \proofstepwidestar[12,13]{x\in A\land y\in A}{\rAI{12,13}}
  \proofstepwidestar[12,13]{(x,y)\in A\times A}{%
    \FormulaRefAuto{(a,b)\in A\times B\eqvdash
      a\in A\land b\in B}[thm]{14}}
  \proofstepwidestar[12,13]{\pi_1(x,y)=x}{%
    \FormulaRefAuto{(x,y)\in A\times B\vdash
      \pi_1(x,y)=x}[thm]{15}}
  \proofstepwidestar[12,13]{\pi_2(x,y)=y}{%
    \FormulaRefAuto{(x,y)\in A\times B\vdash
      \pi_2(x,y)=y}[thm]{15}}
  \proofstepwidestar[1]{%
    (x,y)\in A\times A\rightarrow
    \BinaryOpFunction{A}{\star}(x,y)
      =\pi_1(x,y)\star\pi_2(x,y)}{\rUE{11}}
  \proofstepwidestar[1,12,13]{%
    \BinaryOpFunction{A}{\star}(x,y)
      =\pi_1(x,y)\star\pi_2(x,y)}{\rRE{18,15}}
  \proofstepwidestar[12,13]{%
    \pi_1(x,y)\star\pi_2(x,y)=x\star y}{\rIE{16,17}}
  \proofstepwidestar[1,12,13]{%
    \BinaryOpFunction{A}{\star}(x,y)=x\star y}{\rChain{19,20}}
  \proofstepwidestar[1]{%
    \forall x,y\in A\,
    \BinaryOpFunction{A}{\star}(x,y)=x\star y}{%
    \rUI{\rRI{12,\rUI{\rRI{13,21}}}}}
  \closeproofpartwide
\end{tabproofsplitwide}
